\chapter{PROFIL INDUSTRI}

\section{Profil Politeknik Enjinering Indorama}
Politeknik Enjinering Indorama (PEI) adalah institusi pendidikan tinggi vokasi yang berlokasi di Kabupaten Purwakarta, Jawa Barat. Sesuai dengan tujuan utamanya untuk menjadi politeknik kelas dunia, PEI mengintegrasikan secara langsung pendidikan akademis berkualitas dengan pengalaman praktis di lingkungan industri. Pendekatan ini dirancang untuk menjembatani kesenjangan antara teori rekayasa dan kebutuhan nyata di lapangan kerja industri manufaktur.

Secara spesifik, pelaksanaan Kerja Praktik ini bernaung di bawah Program Studi Teknologi Mekatronika PEI. Program studi jenjang Sarjana Terapan ini menyelenggarakan pendidikan yang berfokus pada disiplin ilmu rekayasa sinergis, menggabungkan teknik mesin, elektronika, sistem kendali, dan rekayasa perangkat lunak. Laboratorium Teknologi Mekatronika PEI memiliki fasilitas tingkat lanjut yang dirancang untuk mereplikasi lingkungan manufaktur modern, menjadikannya tempat yang sangat ideal bagi mahasiswa praktikan untuk mengembangkan purwarupa otomasi dan komputasi cerdas.

\section{Sejarah Politeknik Enjinering Indorama}
PT Indo-Rama Synthetics Tbk. didirikan pada tahun 1975 dan memulai operasi komersial dengan mendirikan pabrik pemintalan kapas di Purwakarta pada tahun 1976. Perusahaan terus melakukan diversifikasi dan memperluas bisnis untuk pasar global. Indorama merupakan salah satu eksportir terbesar di Indonesia dan telah menjadi pemenang rutin penghargaan Primaniyarta yang prestisius untuk kinerja ekspornya, menjadikannya salah satu produsen poliester paling kompetitif di seluruh dunia berkat proses investasi ulang yang berkesinambungan dan program peningkatan produktivitas.

Melalui Yayasan Pendidikan Indorama, PT Indo-Rama Synthetics Tbk. memutuskan untuk mengembangkan politeknik kelas dunia yang berfokus pada penyediaan pendidikan berkualitas tinggi dan relevan dengan industri dengan biaya terjangkau. Lembaga nirlaba ini disebut Politeknik Enjinering Indorama (PEI), yang pendiriannya diselesaikan setelah dua tahun perencanaan cermat mengenai kompetensi pekerjaan dan persyaratan keterampilan industri. PEI didirikan sebagai \textit{eco-campus} di atas lahan hijau seluas 6 hektar dengan lokasi strategis yang berdekatan dengan fasilitas manufaktur Indorama.

Di dalam ekosistem kampus yang terintegrasi dengan industri ini, Program Studi Teknologi Mekatronika dibentuk secara strategis sebagai jawaban atas tantangan era Industri 4.0. Sejak didirikan, program studi ini terus mengembangkan kapasitas laboratoriumnya---mulai dari sistem pneumatik dasar hingga robotika mutakhir---untuk memastikan lulusannya siap menjadi pelopor rekayasa teknologi dan memiliki daya saing di kawasan regional.

\section{Visi dan Misi Politeknik Enjinering Indorama}
Dalam menyelenggarakan pendidikan dan riset terapan, institusi dan program studi menetapkan visi dan misi sebagai arah gerak dan standar keunggulan institusi:

\subsection{Visi Politeknik Enjinering Indorama}
Menjadi Politeknik Kelas Dunia dengan Memberikan Pendidikan Berkualitas dan Pengalaman Praktis di Industri untuk Menghasilkan Lulusan Siap Kerja.

\subsection{Misi Politeknik Enjinering Indorama}
\begin{enumerate}[label=\arabic*.]
    \item Menyelenggarakan pendidikan vokasi untuk menghasilkan sumber daya manusia yang berkualitas, memiliki keunggulan dan kemandirian, dengan didukung pengembangan kurikulum, laboratorium dan Alat Pembelajaran serta Tenaga Pengajar Berkualitas.
    \item Menyelenggarakan penelitian untuk mengembangkan ilmu pengetahuan dan teknologi guna meningkatkan kesejahteraan masyarakat.
    \item Menyelenggarakan pengabdian kepada masyarakat dan diseminasi teknologi dan hasil-hasil penelitian untuk memecahkan masalah-masalah yang ada di masyarakat.
    \item Membangun 5 kompetensi kritikal yang berhubungan dengan penyelesaian masalah, kolaborasi kerjasama tim, produktivitas, fokus dalam kualitas dan memiliki kemampuan komputer dan bahasa Inggris.
\end{enumerate}

\subsection{Visi Program Studi Teknologi Mekatronika}
Menjadi program studi vokasi terdepan yang bermitra erat dengan industri guna mengaplikasikan teknologi di bidang Mekatronika dan memenuhi kebutuhan Sumber Daya Manusia yang kompeten dan mampu yang berdaya saing dalam kawasan regional.

\subsection{Misi Program Studi Teknologi Mekatronika}
\begin{enumerate}[label=\arabic*.]
    \item Menyelenggarakan pendidikan Sarjana Terapan bidang rekayasa mekatronika, mampu menyediakan tenaga profesional yang kompeten di bidang Mekatronika.
    \item Berperan aktif menyelenggarakan penelitian dan kajian terapan dalam bidang rekayasa mekatronika.
    \item Menjalin kemitraan dan bisnis dengan industri tekstil dan industri manufaktur umum melalui rekayasa teknologi mekatronika.
    \item Menciptakan iklim kompetisi bidang rekayasa teknologi yang berwawasan wirausaha.
    \item Mengembangkan teknologi mekatronika yang dapat diaplikasikan secara langsung kepada masyarakat.
\end{enumerate}

\section{Roadmap Politeknik Enjinering Indorama}
Dalam mendukung pencapaian visi institusi dan memfasilitasi pengembangan teknologi tepat guna, Politeknik Enjinering Indorama memiliki peta jalan (\textit{roadmap}) strategis yang komprehensif. Peta jalan ini mencakup tiga pilar utama yang sangat relevan dengan fokus riset Program Studi Teknologi Mekatronika:

\begin{enumerate}[label=\arabic*.]
    \item \textbf{Roadmap Program Pendidikan dan Kurikulum}:
    \begin{itemize}
        \item \textbf{Pendidikan Vokasi Terapan}: Menyelaraskan rancangan kurikulum dengan Standar Nasional Pendidikan Tinggi (SN-Dikti) dan standar kompetensi industri global.
        \item \textbf{Fokus Program Studi}: Berkonsentrasi pada program studi berbasis rekayasa teknik dan manufaktur yang sangat dibutuhkan oleh sektor industri, di antaranya Teknologi Rekayasa Manufaktur, Teknologi Rekayasa Mekatronika, dan Teknologi Listrik.
        \item \textbf{Target Lulusan}: Menciptakan tenaga ahli vokasi yang cerdas, aktif, disiplin, religius, serta siap pakai untuk memimpin di lingkungan kerja industri modern.
    \end{itemize}
    
    \item \textbf{Roadmap Penelitian Terapan}:
    \begin{itemize}
        \item \textbf{Fokus Riset}: Penelitian di PEI diarahkan secara spesifik pada riset terapan yang mampu menjawab tantangan dan kebutuhan nyata di dunia industri, khususnya dalam bidang otomasi dan teknologi manufaktur.
        \item \textbf{Dukungan Institusi}: Kegiatan riset didukung penuh oleh Lembaga Penelitian dan Pengabdian kepada Masyarakat (LPPM) yang mengarahkan pendanaan ke riset yang berpotensi menghasilkan paten dan hilirisasi produk.
        \item \textbf{Kolaborasi Industri}: Terus memperkuat kerja sama strategis dengan mitra korporasi, termasuk PT Indo-Rama Synthetics Tbk, sebagai bagian dari program \textit{Corporate Social Responsibility} (CSR) berbasis pendidikan berkelanjutan.
    \end{itemize}
    
    \item \textbf{Rencana Strategis (Renstra) dan Akreditasi}:
    \begin{itemize}
        \item \textbf{Pencapaian Akreditasi}: Berkomitmen penuh untuk meraih dan mempertahankan akreditasi unggul dari Lembaga Akreditasi Mandiri (LAM) Teknik dan BAN-PT, guna mengukuhkan PEI sebagai institusi rujukan di Jawa Barat.
        \item \textbf{Sistem Penjaminan Mutu}: Melaksanakan proses evaluasi dan audit mutu internal secara berkala melalui Badan Penjaminan Mutu Internal (BPMI).
    \end{itemize}
\end{enumerate}
